\chapter*{Abstract}
\addcontentsline{toc}{chapter}{Abstract}

Il linguaggio Racing Choreographies estende il modello di programmazione coreografica introducendo il costrutto di race, che consente di rappresentare competizione tra processi concorrenti per fornire un valore a un ricevente, mantenendo traccia sia del vincitore sia del perdente e permettendo la successiva sincronizzazione tramite discharge.

Questa tesi sperimentale presenta la progettazione e l’implementazione di un tool per rendere eseguibili programmi in Racing Choreographies. Il lavoro integra la definizione di una grammatica ANTLR4, un parser in C++17 con costruzione di un AST e validazione strutturale post-parsing e un simulatore basato su una macchina astratta che implementa una semantica operazionale semplificata di Racing Choreographies direttamente eseguibile.

